% 第一章：基础概念（基础篇）
% 示范正文章节的常用元素：小节、定理环境、公式、例题与章末习题。
\chapter{基础概念}

\section{什么是示例问题}
% 开头用一小节引出本章要讲的核心概念
本章介绍全书的基础概念。先用一段话给出直观解释，
再给出严格定义，最后通过例子加深理解。

\begin{definition}[示例问题]
  给定样本集合 $\{(\vect{x}_i, y_i)\}_{i=1}^{N}$，
  示例问题要求寻找参数 $\vect{\theta}$，
  使目标函数 $\mathcal{L}(\vect{\theta})$ 最小。
\end{definition}

\section{基本性质}
% 定理 + 证明的标准组合，证明可留作习题
\begin{theorem}[收敛性]
  当学习率 $\eta$ 满足 $\eta < 1/L$（$L$ 为 Lipschitz 常数）时，
  梯度下降法收敛到局部极小点。
\end{theorem}
\begin{proof}
  由凸性与 Lipschitz 连续性，对目标函数应用标准下降引理即得。
\end{proof}

\section{应用举例}
% 例题环境：先给题目再给解答
\begin{example}[一元情形]
  \label{ex:onedim}
  设 $f(x) = x^2$，从 $x_0 = 2$ 出发，学习率 $\eta = 0.25$，
  求一步梯度下降后的位置。
\end{example}
\begin{proof}[解]
  梯度为 $f'(x) = 2x$，于是
  $x_1 = x_0 - \eta f'(x_0) = 2 - 0.25 \times 4 = 1$。
\end{proof}

\section{本章小结}
% 一段话总结 + 引出下一章
本章给出了示例问题的定义与基本收敛性质，
下一章将在更一般的条件下讨论改进算法。

% ---- 章末习题：讲义常用结构，不需要时可删除 ----
\section*{习\quad 题}
\addcontentsline{toc}{section}{习题}
\begin{enumerate}
  \item 证明：当目标函数为强凸函数时，梯度下降线性收敛。
  \item 编程实现一元梯度下降，并复现例~\ref{ex:onedim} 的计算。
\end{enumerate}
